\section{Introduction}\label{sec:intro}

A model deployed in lending, hiring, health care or public administration is rarely judged on accuracy alone. Auditors and regulators also ask whether its errors fall evenly on the groups it affects, and the usual way of putting that question into numbers is to compare group-specific true and false positive rates, as in the equal opportunity and equalised-odds criteria of \citet{hardt2016}. Unlike demographic parity, these criteria depend on the true outcome, so every instance that enters the calculation must be labelled, and labels are often expensive: a clinician may have to review a record, or an investigator to resolve a case. The question also seldom concerns one model. A team typically has several candidate classifiers of similar accuracy, trained with different methods or settings, and wants to know which is fairer, by how much, and whether the difference can be told apart from sampling noise. This paper asks how to answer that with as few labels as possible.

Fairness comparisons need many labels for two reasons. Group rates are ratios with small denominators: a true positive rate in a minority group is estimated from that group's positives only, which may be a few per cent of a test sample. And the quantity of interest is a difference of differences, the gap between groups for one model compared with the gap for another, so the variance of four group rates enters every contrast. Statistical work on fairness metrics has noted the resulting instability, particularly for small groups \citep{lum2022,ferrara2025}. An earlier study of ours of nine models on twelve data settings (the eleven used below plus Nomao, which is excluded later; not part of this paper, results in \texttt{results/realdata\_validation.json}, \texttt{results/realdata\_multi.json} and \texttt{results/certify\_realdata.json}), using held-out pools as a reference population, gave two illustrations. With 500 randomly drawn test labels, the model with the best equal-weight composite of error, calibration, robustness and equalised-odds gap coincided with the reference winner in a median of 36\% of test samples (range 12--96\% across settings). And when the test sample size needed to certify, on a single criterion, that a chosen model is within 0.01 of the best was extrapolated from pilot certificates, the median requirement was about 1,100 random labels for error but about 135,000 for the equalised-odds gap. A comparison that is cheap for accuracy is two orders of magnitude more expensive for fairness under simple random sampling. Budgets of that size are rarely available, so the labels have to be spent where they are informative.

Two facts about the problem make this possible. First, competing models disagree on few instances. For a pair of models $m$ and $j$, the difference in true positive rate within group $g$ is a sum over the positives of that group of $\hat y_{mi}-\hat y_{ji}$, divided by the number of positives $P_g$; the summand is zero wherever the two models make the same prediction. The same holds for false positive rates and for error. Among the three most accurate of our nine models, the average pairwise disagreement rate is between 3.0\% and 9.5\% of the pool in ten of our eleven settings and essentially zero in the eleventh. Near-equally accurate models are known to disagree on individual predictions \citep{marx2020}; here the disagreements are numerous enough to matter for fairness but rare enough to be targeted. Second, group membership is known before labelling. Unlike the outcome, the protected attribute is usually recorded with the instance, so the pool can be divided by group in advance and a small group sampled more heavily than its share. Simple random sampling uses neither fact and spends most of its budget on instances where all models agree and on the majority group.

The problem is as follows. We are given a pool of $N$ unlabelled instances, $M$ fixed classifiers with their binary predictions and predicted probabilities on every instance, and the group of every instance. A budget of $L$ labels may be spent once, before any label is seen. We wish to choose which instances to label and how to estimate, for every pair of models, the difference between their group gaps in true positive rate, the difference between their group gaps in false positive rate, and the difference between their error rates, together with confidence intervals that are valid under the sampling design. Inference is conditional on the trained models, and the target is the value of each contrast on the pool, the natural reference when the pool is the population of use or a large sample from it.

Our method, PATTERN-DIF, has two components. The first is a design. The pool is stratified by group and by the number of models that predict the positive class, so that the instances on which all models agree form the two agreement strata of each group and the strata in between collect the instances on which the models split. Labels are allocated across strata by a weighted generalised Neyman rule that uses the models' averaged predicted probabilities as a proxy for the unknown labels, and are drawn by stratified simple random sampling without replacement. The second component is an estimator. Every contrast is a smooth function of a few population totals, and the proxy means of those totals are known for the whole pool before any label is bought. The stratified difference estimator, in the model-assisted tradition of survey sampling, adds the stratified expansion of the labelled residuals to the proxy total and is design-unbiased for any proxy, however poor. Contrasts are obtained by plug-in and variances by linearisation, which gives Wald intervals. The plain expansion estimator with the same design, PATTERN, is the special case without model assistance. The agreement strata are never left empty: the denominators $P_g$ and $Q_g$ of the group rates involve every instance of the group, including those on which all models agree, so a design that labelled only disagreements would identify the numerators of the contrasts but not the rates themselves.

No ingredient is new on its own, and we state the contribution narrowly. Sampling where models disagree in order to compare their risks goes back at least to \citet{sawade2012active}, who derived the power-optimal sampling distribution for the difference in risk of two models, noted that it favours instances of disagreement, and extended it heuristically to several models; related ideas appear in active model selection \citep{hara2024,kay2025coda} and in discordant-pair evaluation \citep{musgrove2023discordant}. In biostatistics, comparing two screening tests by verifying discordant or test-positive subjects is decades old and covers sensitivity and specificity, as ratios when the doubly negative cell is not verified \citep{schatzkin1987,chock1997,pepe2001comparing}, and optimal two-phase allocation for these quantities has been worked out \citep{mcnamee2002}. Stratification on classifier scores with optimal allocation is established for estimating the performance of one classifier \citep{bennett2010online,kumar2018classifier,fisch2024stratppi}; the use of predictions on a whole pool to sharpen an estimator from a small labelled sample is the idea of prediction-powered inference \citep{angelopoulos2023ppi,fisch2024stratppi} and, much earlier, of model-assisted survey estimation \citep{sarndal1992,breidt2017}; and label-efficient estimation of group true and false positive rates for one model, or for several models with a given labelled set, has been studied with Bayesian, semi-supervised and mixture-model methods \citep{ji2020fairness,gao2025,shanmugam2025}. We add the combination of these ingredients, adapted to ratio-type group-fairness contrasts between several models under a fixed labelling budget, with theory and empirical evidence on its performance. Specifically, we (i) define a stratification that crosses the protected group with the models' agreement pattern and an allocation that minimises, up to finite-population corrections and integer rounding, a weighted sum of the linearised variances of all pairwise contrasts at once; (ii) give design-based theory for the plain and the difference estimators, including the way the linearised variance of a ratio contrast draws on disagreement strata through its numerator and on all strata of a group through its denominator, the condition under which the difference estimator beats the plain one (a within-stratum slope of the influence value on its proxy above one half), the design effect, the allocation, and the effect of a misspecified proxy; (iii) give Wald, non-inferiority and simultaneous max-$t$ tests of the contrasts, with proofs of asymptotic size, a power function whose efficiency ratio equals the design effect (so that the labels needed for a given power shrink by that factor), strong familywise error control over all contrasts, and the allocation that maximises the worst-case power; and (iv) evaluate the method on eleven data settings with nine models, in simulations with a known truth, and against baselines and ablations that isolate where the gain comes from. We do not claim novelty for the use of disagreement, for proxy-based Neyman allocation, or for model-assisted estimation with model predictions, and this is not the first label-efficient method for fairness estimation or for comparing several models.

Most of the empirical gain comes from the design. Efficiency is the ratio of the mean squared error under simple random sampling to that under the method, that is, the factor by which simple random sampling would need more labels; within a setting and budget we take the median over the contrasts that the pool does not determine identically, then the median over settings. For PATTERN-DIF it is 5.8, 6.3, 6.9 and 15.5 at budgets of 250, 500, 1000 and 2000 labels (eleven, ten, eight and six settings; 5.5, 5.8, 6.0 and 9.4 without law school, where the leading models agree on almost every instance), against 3.7, 3.9, 4.4 and 7.3 for the plain estimator on the same design, 1.8 to 2.0 for the difference estimator on a simple random sample, and 1.0 to 1.3 for stratification by group alone with the plain estimator. Stratified prediction-powered inference \citep{fisch2024stratppi} on score strata is competitive with plain PATTERN (higher efficiency in 18 of 35 cells) but PATTERN-DIF beats it in all 35; the importance-sampling design of \citet{sawade2012active} and other alternatives are compared in Section~\ref{sec:baselines}. Efficiency is larger when the models disagree less (Spearman correlation $-0.58$ at $L=250$ without law school). Coverage of nominal 95\% intervals is near nominal on average (0.943 at $L=250$ and 0.949 at $L=2000$), but not for the worst contrast: 0.59 for law school at $L=250$ and 0.82 or more in every other setting at that budget. The evidence has limits: one split and one set of nine models per setting (three further splits for the six largest settings, Section~\ref{sec:splits}); six of the eleven settings use arbitrary median splits as groups, which carry no fairness meaning; the gain varies from 3.5 to 12.1 at $L=250$ outside law school; the variance floor for small strata was designed after seeing the coverage failures, so its benefit is in-sample; and the two-phase result rests on a synthetic population.

The paper is organised as follows. Section~\ref{sec:related} reviews related work and positions the contribution. Section~\ref{sec:setting} fixes the setting, the estimands and the notation. Section~\ref{sec:design} describes the stratification, the allocation and the estimators, and Section~\ref{sec:theory} gives their design-based properties and a numerical check of the theory (Supplementary Section~\SIref{sec:simulation} has the details). Section~\ref{sec:testing} gives the hypothesis tests and their size, power and familywise error properties, and Section~\ref{sec:testing-sim} checks them by simulation. Section~\ref{sec:empirical} reports the study on eleven real data settings, Section~\ref{sec:splits} its repetition over several train--pool splits, and Section~\ref{sec:population} the inference for the deployment population rather than for the pool. Section~\ref{sec:baselines} compares the design with alternative methods from the literature, and Section~\ref{sec:ablations} isolates the contribution of each component. Section~\ref{sec:guidance} turns the results into practical guidance on when the design pays off and how to choose a budget, and Section~\ref{sec:discussion} discusses limitations and open problems. Proofs, additional theory, baseline details, full ablation tables and data details are in the Supplementary Information.
